\documentclass[10pt,a4paper]{amsart}
\usepackage[margin=19mm]{geometry}
\usepackage{amsmath,amssymb,mathtools}
\usepackage[scr=rsfs,cal=euler]{mathalfa}
\usepackage{xfrac}
\usepackage[unicode,hidelinks,bookmarks=false]{hyperref}
\newcommand{\ie}{i.e.\ }
\newcommand{\cf}{cf.\ }
\newcommand{\resp}{respectively\ }
\newcommand{\Subsection}{\subsection}
\providecommand{\nobreakdash}{\nobreak\hbox{-}}
\setlength{\emergencystretch}{2em}
\multlinegap0pt
\usepackage{amssymb}
\newtheorem{lemma}{Lemma}
\def\H{{\bf H}}
\def\HH{{\mathcal H}}
\def\R{{\mathbb R}}
\title{Invariant sets of the double Hilbert transform}
\author{Evgeny Abakumov, Komla Domelevo, Stefanie Petermichl, Alexei Poltoratski}
\date{Source: arXiv:2609.15155v1 (CC BY 4.0)}
\begin{document}
\maketitle

\noindent\emph{Source and licence.} Invariant sets of the double Hilbert transform. Version arXiv:2609.15155v1 (CC BY 4.0), distributed by arXiv under the Creative Commons Attribution 4.0 International licence, http://creativecommons.org/licenses/by/4.0/, as recorded in the arXiv licence metadata for this exact version and shown on the abstract page https://arxiv.org/abs/2609.15155v1. Full licence notice: \texttt{licenses/arxiv-2609.15155-CC-BY-4.0.txt}.\par\smallskip

\noindent\textit{Editorial scope.} Counted unit: the article's lemma L1 (source0.tex lines 104--106). The article's complete original proof of it is delivered with it (source0.tex lines 216--232, one \texttt{proof} environment of 17 lines, which the article places away from the statement). No mathematics is written, repaired or re-derived: the statement and the proof are the article's own lines, in the article's own notation, and no retained line is substituted or re-flowed.  No further source-local context is needed: every cross-reference of every retained line resolves inside the two delivered spans.  Nothing was omitted from any retained span, so the package carries no omission record.  No retained line contains a citation, so the wrapper carries no bibliography and no bibliography entry.  No retained line contains a figure environment, an image-inclusion command or drawing code, so no image is delivered and the package declares no asset dependency.  Editorial formatting only, in addition: an AMS \texttt{amsart} preamble that restates the article's own declarations and macros used by the retained lines, namely the environments \texttt{lemma} on separate counters, together with the standard packages those lines need.  The retained environments print in this document as Lemma 1 in order of appearance, whereas the article prints the same items with the numbering of its own full document.  The complete native source of the article as distributed in the arXiv e-print for this exact version is bundled unchanged as \texttt{sources/210392/source0.tex}.
	\begin{lemma}\label{L1}
		Let $V\subset \R^2$ be a bounded set. Then $\chi_V$ is not an eigenvector of $\H$.
	\end{lemma}

	\begin{proof}[Proof of Lemma \ref{L1}]
		Let $V\subset [-C,C]\times [-C,C]$ and let $l_d=\{(x,y)|y=d\}$ be a horizontal line.
		Notice that for every $d$ such that $|l_d\cap V|>0$ we have
		$$\int_{|x|>C}|\mathcal{H}_1\chi_V(x,d)|^2dx>0.$$
		From the unitarity of the Hilbert transform it follows that
		\begin{align*}
		\iint_{|x|<C}|\mathcal{H}_1\chi_V(x,y)|^2\,dxdy&=\iint_{\R^2}|\mathcal{H}_1\chi_V|^2-
		\iint_{|x|>C}|\mathcal{H}_1\chi_V|^2\\
		&=\|\chi_V\|^2_{L^2(\R^2)}-\iint_{|x|>C}|\mathcal{H}_1\chi_V|^2<\|\chi_V\|^2_{L^2(\R^2)}.
		\end{align*}
		Furthermore, by the unitarity of $\HH_2$ on each vertical line,
		\begin{align*}
		\|\chi_V \H\chi_V\|^2_{L^2(\R^2)}&\leq \|\chi_{\{|x|<C\}} \H\chi_V\|^2_{L^2(\R^2)}\\
		&=\iint_{|x|<C}|\mathcal{H}_1\chi_V(x,y)|^2\,dxdy<\|\chi_V\|^2_{L^2(\R^2)},
		\end{align*}
		which implies that $\H\chi_V\neq \pm \chi_V$.
	\end{proof}

\end{document}
